%! AP Statistics — Unit 9, Lesson 9.1 — Warm-Up (Answer Key)
\documentclass[10pt]{article}
\usepackage{apstats-article}
\usepackage{apstats-key}

\begin{document}

\pageheader{Unit 9, Lesson 9.1}{Warm-Up --- Answer Key}
\namedateperiod

\vspace{0.15in}

A researcher records the number of hours students studied ($x$) and their score
on a statistics exam ($y$) for a random sample of 30 students. The LSRL is:
\[
    \hat{y} = 48.3 + 4.2x
\]

\begin{enumerate}

  \item Identify the slope and $y$-intercept of the LSRL. Interpret the
        \textbf{slope} in context.

        Slope: \ans{4.2} \quad $y$-intercept: \ans{48.3}

        \vspace{0.05in}
        Interpretation: \ans{For each additional hour studied, the predicted exam score
        increases by 4.2 points.}

        \vspace{0.1in}

  \item Use the LSRL to predict the exam score for a student who studied
        for 5 hours. Show your calculation.

        \vspace{0.05in}
        $\hat{y} = 48.3 + 4.2($\ans{5}$) = $ \ans{69.3}

        \vspace{0.1in}

  \item The slope $b = 4.2$ comes from one sample of 30 students. If we took a
        different random sample of 30 students from the same population, would
        we get exactly $b = 4.2$ again? What does this suggest about the
        relationship between the sample slope and the population slope $\beta$?

        \ansline{No --- a different sample would likely produce a different slope.
        This means $b$ is a statistic that varies from sample to sample, just like
        $\bar{x}$ or $\hat{p}$. The true population slope $\beta$ is a fixed (but
        unknown) parameter; $b$ is our sample estimate of it.}

\end{enumerate}

\begin{teachernote}
    Key takeaway: students should leave the warm-up recognizing that $b$ is a
    sample statistic with sampling variability. The central question of Unit 9 ---
    whether $\beta$ could be 0 --- follows naturally from this idea.
\end{teachernote}

\end{document}
